\appendix

\section{Estructura del repositorio y ejecución}\label{anexo:estructura}

\textbf{Cómo ejecutar} (Python 3.10 o superior, sin instalar dependencias):

\begin{verbatim}
python servidor.py                            # abre http://localhost:8000
python herramientas/demo.py                   # plan en consola
python -m unittest discover -s pruebas -v     # pruebas automaticas
python pruebas/evaluacion.py --figuras        # evaluacion del Hito 3
python documento/generar_tablas.py            # tablas y cifras del documento
python herramientas/capturar.py               # capturas de pantalla (Edge o Chrome)
python herramientas/generar_diagramas.py      # diagramas (requiere matplotlib)
\end{verbatim}

\textbf{Estructura:}

\begin{verbatim}
motor/
  semestria.py      planificador: creditos, huecos, entregas, EDF, advertencias
  calibracion.py    factor aprendido de las horas reales
  asesor.py         validacion de horarios y recomendaciones
  asistente.py      chat: intenciones, 12 metodos de estudio, entrevista, modelo local
  progreso.py       estadisticas de progreso y rachas
servidor.py         servidor local (biblioteca estandar) y API JSON
prototipo/          interfaz: index.html, app.js, asistente.js, estilos.css, asistente.css
datos/              horario_ejemplo.json (horario real del estudiante)
pruebas/            test_motor.py, test_mvp.py, test_asistente.py, evaluacion.py
resultados/         evaluacion.json y evaluacion.md (generados)
documento/          este documento (LaTeX), con sus figuras en img/
herramientas/     demo en consola, capturas y diagramas (scripts)
archivo_hito2/      prototipo navegable del Hito 2 (con motor en JavaScript)
\end{verbatim}

\section{Formato del archivo de horario}

El horario es un archivo JSON. Los campos obligatorios son \texttt{semestre\_inicio} y \texttt{bloques}; el resto se completa con valores por defecto al importarlo (la validación informa qué falta).

\begin{verbatim}
{
  "semestre_inicio": "2026-08-24",
  "materias": [ {"codigo": "ISIS ARSW", "creditos": 4} ],
  "bloques": [
    {"materia": "ISIS ARSW", "tipo": "LEC", "dia": 1,
     "inicio": "11:30", "fin": "13:00", "salon": "C C1-203"}
  ],
  "labs": { "ISIS ARSW": {"frecuencia": "semanal", "semanas": null,
                          "dias_entrega": 7, "dificultad": "Dificil"} },
  "entregas_clase": {},
  "disponibilidad": { "0": {"disponible": true, "inicio_trabajo": "08:00"} },
  "trabajos": [ {"id": "m1", "materia": "ISIS ARSW", "tipo": "Parcial",
                 "titulo": "Parcial 2", "entrega": "2026-10-20",
                 "dificultad": "Medio", "horas": 6} ],
  "hechas": [], "en_curso": [], "historial": []
}
\end{verbatim}

Los días se numeran de 0 (lunes) a 6 (domingo); el tipo de bloque es \texttt{LEC} (clase teórica) o \texttt{LAB} (laboratorio); el campo \texttt{horas} de un trabajo es opcional. Los valores de dificultad válidos son «Fácil», «Medio» y «Difícil» (con tilde).
